\documentclass[12pt]{article}

% Packages
\usepackage[
  paperwidth=11in,
  paperheight=8.5in,
  margin=0.5in
]{geometry}
\usepackage{amsmath, amssymb}
\usepackage{booktabs}
\usepackage{tabularx}

\newcommand{\tacred}{FS-TACRED}
\newcommand{\qwensmall}{Qwen3-4B}
\newcommand{\gemmasmall}{Gemma3-4B}

\begin{document}


\section{Experiment recep} I expored 4 strategies to optimize a prompt:
\begin{enumerate}
    \item Update using the feedback on corrects and mistakes
    \item Update using the feedback on corrects 
    \item Update using the feedback on mistakes
    \item Update by random changes in the prompt (no feedback on corrects mistakes) 
\end{enumerate}

\section{Analysis on Results} 
\begin{itemize}
    \item Just applying some random changes is the most efficient approach. The optimzed prompts are quite very similar to the initial prompt with minor changes in the statements.
    \item Using feedback on both corrects and mistakes was 2nd best approach. The optimzed prompts are quite larger than the initial prompt with many instructions (which were learnt from the feedbacks).
    \item Using feedback from corrects was least efficient approach.
\end{itemize}

\section{Discussions and next works}
\begin{itemize}
    \item We has started from a standard baseline prompt i.e. it already has basic details (e.g. what is the task) and instruction (e.g. how to decide yes / no).
    The optmized prompt doesn't change much from this baseline prompt but mostly contains only additional instructions.
    However, the GEPA paper starts from a simplest prompt that mostly contain one or two sentences. Then, the optimized prompts becomes larger with many instructions (It can evolve from the baseline in a free way).
    I was thinking should we also try a simplest prompt (an example is shown in Table~\ref{tab:example-initial-prompt})
\end{itemize}


\section{Prompts}
\label{sec:prompts}

\begin{table}[h]
\centering
\begin{tabular}{ll}
\toprule
\textbf{Type} & \textbf{Text} \\
\midrule
Given a relation, a support example, and a query, classify whether the relation holds in the query sentence. Answer with yes or no only. \\
\bottomrule
\end{tabular}
\caption{Example initial prompt similar to GEPA paper}
\label{tab:example-initial-prompt}
\end{table}

\paragraph{Notes on GEPA paper}
\begin{itemize}
    \item They start from a simlpest instruction (i.e. few sentence only) and their prompt results into a large prompt after optimization.
    That is why I see that the scores of the child nodes are significantly higher than the root node.
    \item The best resulting prompt may come from any path .. So an extensive search (exploration+eploitation) of the tree is highly important
    \item ** Bascially, if a actual prompt has two parts: Instruction and Input, they optimize the Intruction part of the prompt. ***
\end{itemize}



% \bibliographystyle{plain}  
% \bibliography{custom}    

\end{document}
